\documentclass[tikz,border=6pt]{standalone}
\usepackage{peabra}
\begin{document}
\begin{tikzpicture}[x=1cm,y=1cm,
  gy/.style={sotetdoboz=30mm, font=\small\bfseries},
  ag/.style={doboz=27mm, font=\footnotesize\bfseries, minimum height=7mm},
  lv/.style={font=\scriptsize, align=left, anchor=north west, text width=44mm, inner sep=1pt}]
  \node[gy] (g) at (0,0) {Üzembiztonság\\{\scriptsize\mdseries (dependability)}};
  \node[ag] (a) at (-6.2,-1.6) {Tulajdonságok};
  \node[ag] (f) at (0,-1.6) {Fenyegetések};
  \node[ag] (e) at (6.2,-1.6) {Eszközök};
  \foreach \x in {a,f,e} \draw[nyilv] (g.south) -- (\x.north);
  \node[lv] at (-8.4,-2.15) {\textbf{rendelkezésre állás} -- kész a működésre\\
    \textbf{megbízhatóság} -- folyamatos helyes működés\\
    \textbf{biztonság} -- nincs katasztrofális következmény\\
    \textbf{integritás} -- nincs nem megengedett változás\\
    \textbf{karbantarthatóság} -- javítható, módosítható};
  \node[lv] at (-2.2,-2.15) {\textbf{hibaok} (fault)\\ \hspace*{3mm}$\downarrow$ aktiválódik\\
    \textbf{hiba} (error) -- hibás belső állapot\\ \hspace*{3mm}$\downarrow$ eléri a kimenetet\\
    \textbf{meghibásodás} (failure) -- a szolgáltatás eltér a helyestől};
  \node[lv] at (4.0,-2.15) {\textbf{hibamegelőzés} -- gondos fejlesztés\\
    \textbf{hibatűrés} -- redundancia, diagnosztika\\
    \textbf{hibaeltávolítás} -- verifikáció, tesztelés\\
    \textbf{hibaelőrejelzés} -- megbízhatósági modellek, elemzés};
\end{tikzpicture}
\end{document}
